\documentclass[12pt,a4paper]{article}
\usepackage[utf8]{inputenc}
\usepackage[portuguese]{babel}
\usepackage[T1]{fontenc}
\usepackage{amsmath}
\usepackage{amsfonts}
\usepackage{amssymb}
\usepackage{graphicx}
\usepackage{booktabs}
\usepackage{hyperref}
\usepackage{geometry}
\geometry{margin=2.5cm}

\title{\textbf{Relatório sobre consumo de APIs}}
\author{}
\date{}

\begin{document}

\maketitle

\section{IBGE}

Essa API basicamente serve para extrair dados categóricos e numéricos sobre varias coisas relacionadas ao Brasil. Nessa API é possível pesquisar coisas como

\subsection{Censos e Contagens Populacionais}
\textit{Censo Demográfico, Contagem da População, Censo Comum do Mercosul, Estimativas e Projeção da População}

\textbf{Dados:} tamanho populacional, idade, renda, escolaridade, moradia por região/município.

\textbf{Uso:} planejamento urbano, repasses federais, estudos de mercado, projeção de demanda.

\subsection{Índices de Preços}
\textit{INPC, INPC-E, IPCA, IPCA-15, IPCA-E, Índice de Reajuste do Salário Mínimo, IPP, Preços ao Consumidor em Real, SINAPI}

\textbf{Dados:} variação de preços ao consumidor, produtor e construção civil.

\textbf{Uso:} correção de contratos, reajuste salarial, política monetária, precificação.

\subsection{Contas Nacionais e PIB}
\textit{Contas Nacionais Anuais, Trimestrais, PIB dos Municípios}

\textbf{Dados:} crescimento econômico, PIB setorial, consumo, investimento.

\textbf{Uso:} análise macroeconômica, decisões de investimento, comparação regional.

\subsection{Pesquisas Industriais}
\textit{PIA, PIA-Empresa, PIA-Produto, PIM Produção Física, PIM Dados Gerais, PIM Emprego e Salário}

\textbf{Dados:} produção, receita, custos, emprego e salários industriais (mensal e anual).

\textbf{Uso:} análise de competitividade, tendências setoriais, planejamento de produção.

\subsection{Comércio e Serviços}
\textit{PAC, PAS, PMC, PMS, Pesquisas de Hospedagem, Publicidade, TI}

\textbf{Dados:} receita, pessoal ocupado, custos do comércio e serviços.

\textbf{Uso:} dimensionamento de mercado, análise de consumo, estratégias comerciais.

\subsection{Produção Agrícola}
\textit{PAM, LSPA, Produção de Ovos, Extração Vegetal e Silvicultura}

\textbf{Dados:} área plantada, produtividade, safras por cultura e município.

\textbf{Uso:} planejamento de safras, precificação agrícola, seguro rural.

\subsection{Produção Pecuária}
\textit{PPM, Abate de Animais (trimestral/mensal), Leite (trimestral/mensal), Couro}

\textbf{Dados:} rebanhos, abate, produção de leite e couro.

\textbf{Uso:} planejamento frigorífico, formação de preços de carne/leite, políticas pecuárias.

\subsection{Censo e Estudos Agropecuários}
\textit{Censo Agropecuário, Contribuição dos Polinizadores}

\textbf{Dados:} estrutura fundiária, tecnologias agrícolas, dependência de polinização.

\textbf{Uso:} políticas agrárias, extensão rural, valoração ambiental.

\subsection{Contas Ambientais}
\textit{Contas de Ecossistemas, Água, Terra, Energia (biomassa), Espécies Ameaçadas}

\textbf{Dados:} uso de recursos naturais, impactos ambientais das atividades econômicas.

\textbf{Uso:} sustentabilidade, políticas de conservação, economia verde.

\subsection{Saúde}
\textit{PNS, PNS Escolar, Assistência Médico-Sanitária}

\textbf{Dados:} doenças, hábitos de saúde, infraestrutura hospitalar.

\textbf{Uso:} políticas de saúde pública, planejamento hospitalar.

\subsection{Pesquisas Domiciliares e Emprego}
\textit{PNAD, PNAD Contínua (mensal/trimestral/anual), POF, PME}

\textbf{Dados:} emprego, renda, desemprego, gastos familiares.

\textbf{Uso:} políticas trabalhistas, análise de mercado consumidor, segmentação.

\subsection{Empresas e Empreendedorismo}
\textit{CEMPRE, Demografia das Empresas, MEI, FASFIL}

\textbf{Dados:} nascimento/morte de empresas, perfil empresarial, ONGs.

\textbf{Uso:} análise de ambiente de negócios, políticas de fomento.

\subsection{Pesquisas Municipais e Infraestrutura}
\textit{MUNIC, Saneamento Básico, Estatísticas do Registro Civil}

\textbf{Dados:} gestão municipal, saneamento, nascimentos, casamentos, óbitos.

\textbf{Uso:} políticas públicas locais, planejamento de infraestrutura.

\subsection{Inovação e Estoques}
\textit{PINTEC, Pesquisa de Estoques}

\textbf{Dados:} inovação tecnológica empresarial, níveis de estoque.

\textbf{Uso:} políticas de P\&D, gestão de cadeia de suprimentos.

\subsection{Construção Civil}
\textit{PAIC, SINAPI}

\textbf{Dados:} receita, custos, emprego na construção, índices de custos.

\textbf{Uso:} análise do setor imobiliário, orçamentos de obras.

\subsection{Sustentabilidade}
\textit{IDS, ODS, Áreas Urbanizadas}

\textbf{Dados:} indicadores socioambientais, avanço nas metas ONU, expansão urbana.

\textbf{Uso:} monitoramento de desenvolvimento sustentável, planejamento territorial.





\end{document}